\scanpage{191}

% B01 / text / top, heading and first line
\mspage{98}\uline{2.}

Soit \deleted{$(S,A_0)$ un graphe}

% B02 / display / upper, (1)
\[
 (S,A_0)\qquad A_0\subset\unsure{P|\mathcal P|\mathfrak P}_2(S) \tag*{(1)}
\]

% B03 / text / upper, before (2)
un graphe, satisfaisant la condition de ``triangularité''

% B04 / display / upper, (2), left brace
\[
\left\{\begin{aligned}
 &\text{a)}\quad S=\bigcup_{a\in A_0}a\\
 &\text{b)}\quad \forall\,\underset{\substack{\shortparallel\\\{s,t\}}}{a}\in A_0,\quad \exists\ \text{exactement deux éléments }u\text{ de }S,\\
 &\hspace{12em}\text{tels que }\{s,u\},\{t,u\}\in A_0
\end{aligned}\right. \tag*{(2)}
\]

% B05 / text / upper middle, before (3)
\deleted{\unsure{avec}} de telle façon que pour tout $s\in S$, son étoile pour $A_0$

% B06 / display / upper middle, (3)
\[
 \mathrm{Et}_{A_0}(s)=\{t\in S\mid\{s,t\}\in A_0\}\qquad =E \tag*{(3)}
\]

% B07 / text / middle, after (3)
est de façon canonique munie d'une structure de contour \deleted{\unsure{dont}} \unsure{indenté|induit|induite}, dont l'ens des arêtes est

% B08 / display / middle, (4)
\[
 A_{0E}=P_2(E)\cap A_0 \tag*{(4)}
\]

% B09 / text / middle, bracketed passage
dont les composantes connexes sont d'ordre $\geq 3$. [Je suppose donné de plus un groupe d'automorphismes
$G_0$ de $(S,A_0)$, transitif \deleted{sur} $\overrightarrow{A_0}$, et \deleted{simplement}\inserted{même} transitif sur l'ens des répères triangulaires
% slanted
\margin{\unsure{descend|devient|demandé} plus bas}

% B10 / display / lower middle, end of bracket, (5)
\[
 \mathrm{Rep}(\unsure{s|S},A_0)\overset{\mathrm{def}}{=}
 \{(s,t,u)\in S^3\mid \{s,t\},\{t,u\},\deleted{\}}\{u,s\}\in A_0\}\,] \tag*{(5)}
\]

% B11 / text / lower middle, restart
\deleted{Je suppose de plus \unsure{satisf.|satisfaite}} Considérons l'application

% B12 / diagram / lower, (6)
\diagram{(6)\qquad\begin{tikzcd}[ampersand replacement=\&,row sep=0pt]
 {\mathrm{Et}_{A_0}:}\& S\arrow[r,dash]\&P(S)\\
 \&s\arrow[r,mapsto]\&\mathrm{Et}_{A_0}(s)
\end{tikzcd}}

% B13 / text / lower, before (7)
et soit $\unsure{\mathcal E|\mathscr E|\Sigma}_A$ son image

% B14 / display / lower, (7), vertical inclusion
\[
 \underset{\substack{\cap\\P(S)}}{\mathcal E_A}=\{E\in P(S)\mid\exists\,s\in S\ \text{t.q.}\ E=\mathrm{Et}_{A_0}(s)\} \tag*{(7)}
\]

% B15 / text / foot, continuing
On note que pour tout $E\in\mathcal E_A$, la structure de contour
sur $E$ définie par (4) ne dépend pas du choix du $s\in S$ tel que $E=\mathrm{Et}_{A_0}(s)$. On récupère $A_0$, à

% archive sheet number: 190
